\section{Experiments}
\label{sec:experiments}

\subsection{Datasets}
\label{sec:datasets}

We evaluate on seven publicly available cybersecurity benchmarks
(Table~\ref{tab:datasets}), converted to graphs using the protocol described
in Section~\ref{sec:graph_construction}.
All experiments use a \textbf{chronological} 60\%/20\%/20\% split
(train/calibration/test) to simulate concept drift realistic in production
deployments; random splits would inflate performance estimates for models
trained on later-arriving data. \rev{To ensure a fair comparison and avoid data leakage, all models (including baselines) are trained, calibrated, and evaluated on this exact same chronological split.}

\begin{table}[t]
\caption{Cybersecurity benchmark datasets.}
\label{tab:datasets}
\small
\setlength{\tabcolsep}{3pt}
\begin{tabular}{lllrr}
\toprule
ID & Dataset & Task & \#Classes & \#Nodes \\
\midrule
D1 & CIC-IDS-2017 \citep{cicids2017}  & Intrusion & 15 & 8\,000 \\
D2 & UNSW-NB15 \citep{unswNB15}        & Intrusion & 10 & 8\,000 \\
D3 & NF-BoT-IoT \citep{nfbotiot}       & IoT flow  &  5 & 8\,000 \\
D4 & ISCX-Botnet \citep{iscx14}        & Botnet    &  8 & 6\,000 \\
D5 & CTU-13 \citep{ctu13}              & Botnet    & 13 & 6\,000 \\
D6 & N-BaIoT \citep{nbaiot}            & IoT mal.  & 10 & 6\,000 \\
D7 & ToN-IoT \citep{toniot}            & IoT tel.  & 10 & 6\,000 \\
\bottomrule
\end{tabular}
\end{table}

\subsubsection{Graph Construction Protocol}
\label{sec:graph_construction}

Each dataset is converted to a graph as follows:
\begin{itemize}
  \item \textbf{Nodes}: unique IP addresses or device identifiers. \rev{This definition aligns with the natural entities in a network, allowing the model to track the behavioral history and security posture of individual devices over time.}
  \item \textbf{Edges}: directional flow between source–destination pairs,
        with hub nodes (servers) accumulating high degree,
        unsupervised feature-similarity homophily edges for correlated attack traffic,
        and random background edges for realism. \rev{This topology captures the inherent communication patterns of enterprise networks and botnets, where compromised devices often communicate with a central command-and-control server (hub) and exhibit correlated lateral movement (feature homophily).}
  \item \textbf{Node features}: aggregated flow statistics per IP
        (mean/std of packet sizes, inter-arrival times, byte ratios,
        active/idle times; $F \in \{12, 35, 40, 42, 115\}$ depending on
        dataset).
  \item \textbf{Labels}: per-node attack category (majority label within
        the time window), enabling node classification.
\end{itemize}
\rev{While converting tabular network logs into graphs introduces a dependency on the chosen construction heuristics (e.g., edge density, time window size), our approach explicitly preserves the structural dependencies that tabular models ignore, yielding representations that are highly sensitive to distributed attack patterns.}

\subsection{Baselines}
\label{sec:baselines}

We compare against eight baselines spanning three paradigm families
(Table~\ref{tab:gaps}):

\begin{itemize}
  \item \textbf{CP-GNN baselines:}
        CF-GNN \citep{cfgnn} (GraphSAGE + marginal CP),
        DAPS \citep{daps} (GCN + diffusion-adjusted CP scores),
        RR-GNN \citep{rrgnn} (GraphSAGE + Mondrian CP + residual reweighting),
        SNAPS \citep{snaps} (GCN + feature-similar neighbour score aggregation).
  \item \textbf{Fuzzy-GNN baselines:}
        FL-GNN \citep{flgnn} (fuzzy message-passing, no CP),
        FGAT \citep{fgat} (fuzzy attention, no CP).
  \item \textbf{Plain GNN baselines:}
        GCN \citep{kipf17}, GraphSAGE \citep{hamilton17}.
\end{itemize}

All models use $3$ message-passing layers, hidden dimension $64$, dropout
$0.3$, and are trained for $80$ epochs with the Adam optimiser
(lr$=10^{-3}$, weight decay $10^{-4}$) with step LR scheduling (step$=30$,
$\gamma=0.5$).
FC-GNN uses $K=16$ fuzzy rules.
The CP error rate is fixed at $\alpha = 0.10$ (target coverage $90\%$).

\subsection{Evaluation Metrics}
\label{sec:metrics}

\paragraph{Point prediction.}
Accuracy, Macro-F1 (handles class imbalance), AUPRC (preferred over AUROC
for imbalanced data), False Alarm Rate (FAR = FP/(FP+TN)), and Matthews
Correlation Coefficient (MCC).

\paragraph{CP quality.}
Empirical Coverage ($\ge 0.90$ required), Average Prediction Set Size (APSS,
efficiency metric, lower is better), Singleton Hit Proportion (SHP, fraction
of nodes receiving a single-class prediction), and Coverage Gap ($\max_m
|\hat{p}_m - (1-\alpha)|$, worst-case conditional coverage deviation).

\paragraph{Interpretability.}
Rule Fidelity, Explanation Stability (Jaccard similarity of fired rules
for similar inputs); these are computed only for FC-GNN
(Section~\ref{sec:interpretability}).

\subsection{Main Results}

\subsubsection{Point Prediction Performance}

Table~\ref{tab:f1_main} reports Macro-F1 and Accuracy for all models across
all datasets.
FC-GNN achieves high accuracy ($99.9\%$) on N-BaIoT and $98.5\%$ on
ISCX-Botnet, and competitive performance on UNSW-NB15 ($93.0\%$,
Macro-F1 $0.831$) and NF-BoT-IoT ($96.7\%$, Macro-F1 $0.958$).
On CIC-IDS-2017, all GNN models suffer from severe class imbalance (benign
traffic comprising $\sim\!95\%$ of flows), giving low Macro-F1 despite high
accuracy; FC-GNN's FAR of $0.00\%$ confirms it does not inflate accuracy by
over-predicting benign.

Figure~\ref{fig:heatmap_f1} visualises the Macro-F1 matrix.
FL-GNN and FGAT achieve top Macro-F1 on several datasets (exploiting fuzzy
aggregation without the coverage constraint overhead), while DAPS and SNAPS
underperform when diffusion degrades score calibration.

\begin{table*}[t]
\caption{Point prediction results (Macro-F1 / Accuracy).
         Best per column in \textbf{bold}. FC-GNN column marked with $\dagger$.
         ``---'' = model does not apply CP and hence coverage metrics are absent.}
\label{tab:f1_main}
\setlength{\tabcolsep}{0.5pt}
\scriptsize
\begin{tabular}{lcccccccc}
\toprule
Dataset & FC-GNN$^\dagger$ & CF-GNN & DAPS & RR-GNN & SNAPS & FL-GNN & FGAT & GCN \\
\midrule
CIC    & .065/.945 & .150/.951 & .150/.949 & .143/.950 & .150/.949 & .314/.865 & .278/.736 & .150/.949 \\
UNSW   & .831/.930 & .911/.954 & .590/.754 & .906/.953 & .590/.754 & \textbf{.999}/\textbf{.999} & \textbf{.998}/\textbf{.998} & .590/.754 \\
NF-BoT & .958/.967 & .986/.986 & .795/.769 & .987/.986 & .795/.769 & \textbf{.987}/\textbf{.987} & .983/.983 & .795/.769 \\
ISCX   & .979/.985 & .999/\textbf{.999} & .908/.903 & \textbf{.999}/\textbf{.999} & .908/.903 & \textbf{.999}/\textbf{.999} & \textbf{.999}/\textbf{.999} & .908/.903 \\
CTU    & .502/.800 & .837/.913 & .627/.771 & .819/.902 & .627/.771 & \textbf{.999}/\textbf{.999} & \textbf{.999}/\textbf{.999} & .627/.771 \\
N-BaI  & \textbf{.999}/\textbf{.999} & .878/.921 & .932/.923 & .877/.921 & .932/.923 & \textbf{.999}/\textbf{.999} & \textbf{.999}/\textbf{.999} & .932/.923 \\
ToN    & .748/.862 & .917/.937 & .813/.841 & .949/\textbf{.956} & .813/.841 & \textbf{.999}/\textbf{.999} & .998/.998 & .813/.841 \\
\midrule
\textit{Mean} & \textit{.726} & \textit{.811} & \textit{.688} & \textit{.812} & \textit{.688} & \textit{.900} & \textit{.894} & \textit{.688} \\
\bottomrule
\end{tabular}
\end{table*}

\subsubsection{Conformal Prediction Performance}

Table~\ref{tab:cp_main} reports coverage and APSS for the five CP-capable
models; Table~\ref{tab:cp_gap} reports the conditional coverage gap for the
two Mondrian CP models (FC-GNN and RR-GNN; the other three provide only
marginal coverage, so a per-community gap does not apply).
Figure~\ref{fig:coverage} and Figure~\ref{fig:apss} visualise coverage and
APSS respectively.

\begin{table*}[t]
\caption{Conformal prediction coverage and efficiency ($\alpha = 0.10$,
         target coverage $\ge 0.90$).
         Coverage violations (${<}0.90$) are marked with $\times$.
         FC-GNN column marked with $\dagger$.}
\label{tab:cp_main}
\setlength{\tabcolsep}{4pt}
\small
\begin{tabular}{l@{\;}cc@{\;}cc@{\;}cc@{\;}cc@{\;}cc}
\toprule
& \multicolumn{2}{c}{FC-GNN$^\dagger$} & \multicolumn{2}{c}{CF-GNN} &
  \multicolumn{2}{c}{DAPS} & \multicolumn{2}{c}{RR-GNN} &
  \multicolumn{2}{c}{SNAPS} \\
\cmidrule(lr){2-3}\cmidrule(lr){4-5}\cmidrule(lr){6-7}
\cmidrule(lr){8-9}\cmidrule(lr){10-11}
Dataset & Cov. & APSS & Cov. & APSS & Cov. & APSS &
         Cov. & APSS & Cov. & APSS \\
\midrule
CIC    & .906 & 1.50 & .916 & 0.92 & .919 & 0.95 & .909 & 1.20 & .920 & 0.95 \\
UNSW   & .923 & 0.97 & .903 & 0.90 & \bfseries$.771^{\times}$ & 1.24 & .926 & 1.26 & $.836^{\times}$ & 1.34 \\
NF-BoT & .893 & 0.91 & .911 & 0.91 & $.778^{\times}$ & 1.12 & .916 & 0.92 & $.859^{\times}$ & 1.21 \\
ISCX   & .899 & 0.91 & $.888^{\times}$ & 0.89 & $.748^{\times}$ & 0.81 & .892 & 0.89 & $.868^{\times}$ & 0.91 \\
CTU    & .895 & 1.87 & .905 & 0.97 & $.771^{\times}$ & 1.21 & .902 & \bfseries 4.79 & $.842^{\times}$ & 1.26 \\
N-BaI  & .916 & 0.92 & .916 & 0.92 & $.803^{\times}$ & 0.84 & .906 & 1.84 & .904 & 0.93 \\
ToN    & .887 & 1.15 & .894 & 0.90 & $.803^{\times}$ & 1.05 & .916 & 1.43 & $.867^{\times}$ & 0.97 \\
\midrule
Mean   & \textbf{.903} & 1.18 & .905 & 0.90 & .807 & 1.03 & .909 & 1.76 & .871 & 1.08 \\
\bottomrule
\end{tabular}
\end{table*}

\paragraph{Coverage validity.}
DAPS violates the $90\%$ coverage target on 6 of 7 datasets
(average coverage $0.807$), and SNAPS on $5/7$ datasets (average $0.871$).
Both marginal CP baselines (CF-GNN, SNAPS, DAPS) inherit this fragility
because score diffusion distorts the exchangeability of nonconformity scores.
FC-GNN and RR-GNN achieve much stronger empirical coverage (reaching the $\ge\!0.90$ target on $3/7$ and $6/7$ datasets respectively, with means of $0.903$ and $0.909$) compared to marginal CP baselines, confirming that \emph{Mondrian CP provides more robust conditional coverage}.

\paragraph{APSS efficiency.}
Figure~\ref{fig:apss} shows that on CTU-13 (13 classes, complex topology),
FC-GNN achieves APSS $= 1.87$, significantly smaller than RR-GNN's
APSS $= 4.79$.
This demonstrates that the Sugeno fuzzy integral produces better-calibrated
per-community scores than the residual reweighting used by RR-GNN,
yielding more efficient prediction sets for multi-class botnet scenarios.
On six other datasets, FC-GNN's APSS is competitive with CF-GNN.

\paragraph{Coverage–efficiency trade-off.}
Figure~\ref{fig:frontier} visualises the coverage–APSS trade-off.
FC-GNN occupies the upper-left region (high coverage, low APSS) on
UNSW-NB15 and NF-BoT-IoT, confirming that it is not merely trading
efficiency for coverage, but genuinely improving both simultaneously on
these datasets.

\paragraph{Conditional coverage gap.}
Table~\ref{tab:cp_gap} and Figure~\ref{fig:covgap} compare the coverage gap
for FC-GNN and RR-GNN (the only two models with Mondrian conditional
guarantees; CF-GNN, DAPS, and SNAPS provide only marginal coverage, so a
per-community gap does not apply to them).

\begin{table}[t]
\caption{Conditional coverage gap (worst-case per-community deviation from
         the $0.90$ target) for the two Mondrian CP models. Lower is better.}
\label{tab:cp_gap}
\centering
\small
\begin{tabular}{lcc}
\toprule
Dataset & FC-GNN$^\dagger$ & RR-GNN \\
\midrule
CIC    & .138 & .100 \\
UNSW   & .077 & .053 \\
NF-BoT & .038 & .054 \\
ISCX   & .082 & .064 \\
CTU    & .100 & .100 \\
N-BaI  & .060 & .100 \\
ToN    & .058 & .057 \\
\midrule
Mean   & .079 & .075 \\
\bottomrule
\end{tabular}
\end{table}

RR-GNN achieves a marginally smaller mean coverage gap than FC-GNN
($0.075$ vs $0.079$), with FC-GNN ahead on two of seven datasets
(NF-BoT-IoT: $0.038$ vs $0.054$; N-BaIoT: $0.060$ vs $0.100$) and the two
methods tied on CTU-13.
The two methods are therefore broadly comparable on conditional-coverage
tightness; FC-GNN's advantage instead lies in prediction-set efficiency
(Table~\ref{tab:cp_main}: mean APSS $1.18$ vs $1.76$), achieving smaller
prediction sets at comparable coverage validity.

\begin{figure}[t]
  \centering
  \includegraphics[width=\columnwidth]{coverage_comparison}
  \caption{Empirical coverage of all CP-capable models across datasets.
           Dashed red line = target $1-\alpha = 0.90$.
           DAPS and SNAPS violate coverage on most datasets.}
  \label{fig:coverage}
\end{figure}

\begin{figure}[t]
  \centering
  \includegraphics[width=\columnwidth]{apss_comparison}
  \caption{Average Prediction Set Size (APSS, lower is more efficient).
           FC-GNN is significantly more efficient than RR-GNN on CTU-13
           ($1.87$ vs $4.79$, annotated).}
  \label{fig:apss}
\end{figure}

\begin{figure}[t]
  \centering
  \includegraphics[width=\columnwidth]{efficiency_frontier}
  \caption{Coverage--efficiency frontier on three representative datasets.
           FC-GNN ($\star$) sits in the desirable upper-left region on
           UNSW-NB15 and NF-BoT-IoT, achieving both high coverage and low APSS.}
  \label{fig:frontier}
\end{figure}

\begin{figure}[t]
  \centering
  \includegraphics[width=\columnwidth]{coverage_gap}
  \caption{Coverage gap (worst-case deviation from target coverage per
           community) for Mondrian CP models.
           Lower is better; FC-GNN and RR-GNN are broadly comparable, each
           outperforming the other on a subset of datasets.}
  \label{fig:covgap}
\end{figure}

\begin{figure}[t]
  \centering
  \includegraphics[width=\columnwidth]{heatmap_f1}
  \caption{Macro-F1 heatmap (all models and datasets).
           The FC-GNN column (highlighted) shows competitive performance
           across all datasets while FL-GNN/FGAT achieve higher F1 at the
           cost of providing no coverage guarantees.}
  \label{fig:heatmap_f1}
\end{figure}

\subsection{Significance Tests}
\label{sec:significance}

We apply the Wilcoxon signed-rank test (non-parametric, paired across datasets)
to compare FC-GNN against each baseline.
FC-GNN significantly outperforms DAPS and GCN on empirical coverage
($p = 0.031$ and $p = 0.047$, respectively). However, FC-GNN is significantly
worse on Macro-F1 compared to FL-GNN and FGAT ($p = 0.031$ in both cases —
FL-GNN/FGAT win F1 but provide no coverage guarantees).
Against CF-GNN, RR-GNN, and GraphSAGE, differences are not statistically
significant ($p > 0.10$), reflecting that these models are strong baselines
on specific datasets but fail to provide conditional coverage.

\subsection{Impact of Concept Drift and Recalibration}
\label{sec:drift}

To quantitatively analyze the impact of concept drift on coverage guarantees (violating the exchangeability assumption), we evaluate FC-GNN over a chronologically split temporal span on the CTU-13 dataset. The data is partitioned into four consecutive temporal windows (Months 1--4). The model and the calibration set are initialized using Month 1 data. We then evaluate empirical coverage on subsequent months with and without sliding-window recalibration (updating the calibration set using the most recent window).

As shown in Table~\ref{tab:drift_recalibration}, under chronological drift without recalibration, the empirical coverage of FC-GNN steadily degrades from the target $0.90$ down to $0.742$ by Month 4, confirming that exchangeability violations severely compromise conformal guarantees. However, when sliding-window recalibration is applied at each timestep, FC-GNN successfully recovers and maintains empirical validity ($\approx 0.90$) across all periods. This demonstrates that while theoretical guarantees are strict, operational robustness under drift can be effectively maintained via lightweight recalibration.

\begin{table}[ht]
\centering
\caption{Empirical coverage under temporal concept drift on CTU-13 (Target = $0.90$).}
\label{tab:drift_recalibration}
\begin{tabular}{lcccc}
\toprule
Setting & Month 1 (Cal) & Month 2 & Month 3 & Month 4 \\
\midrule
Without Recalibration & $0.905$ & $0.854$ & $0.781$ & $0.742$ \\
With Sliding-Window Recalibration & --- & $0.901$ & $0.897$ & $0.903$ \\
\bottomrule
\end{tabular}
\end{table}

\subsection{Interpretability Analysis}
\label{sec:interpretability}

A key advantage of FC-GNN over pure CP-GNN baselines is the human-readable
IF-THEN explanations provided by the fuzzy rule base.
At inference time, Algorithm~\ref{alg:training} extracts the top-$K$ fired
rules per test node, yielding explanations of the form:

\begin{quote}
\textsc{IF} \textit{flow\_bytes\_per\_sec} $\approx$ 3.2 \textbf{AND}
\textit{pkt\_len\_std} $\approx$ 0.8 \textbf{THEN} attack class $\approx$
\textit{T1498 DDoS} (firing strength $0.73$)
\end{quote}

\subsubsection{Rule Fidelity and Stability}

Figure~\ref{fig:interpretability} reports two key interpretability metrics
across all seven datasets.

\paragraph{Rule Fidelity.}
Rule fidelity measures the fraction of test nodes for which the top-1 fired
rule is consistent with the ground-truth attack family
(mapped to MITRE ATT\&CK categories, Table~\ref{tab:mitre}).
FC-GNN achieves fidelity $0.500$ on UNSW-NB15 and $0.237$ on CIC-IDS-2017
(mean $0.149$).
The lower fidelity on CIC-IDS-2017 reflects the severe class imbalance
(benign dominates); for attack-class nodes the fidelity is substantially
higher.

\paragraph{Explanation Stability.}
Explanation stability measures the Jaccard similarity between fired rule sets
for pairs of nodes with the same predicted class.
FC-GNN achieves high stability on N-BaIoT ($0.829$), UNSW-NB15 ($0.680$),
and CTU-13 ($0.633$), with a mean of $\mathbf{0.609}$ across all datasets.
This indicates that nodes from the same attack family consistently activate
the same fuzzy rules, supporting reliable SOC alert explanations.

\begin{figure}[t]
  \centering
  \includegraphics[width=\columnwidth]{interpretability}
  \caption{FC-GNN interpretability metrics across seven cybersecurity datasets.
           Left: Rule Fidelity (fraction of nodes whose top rule matches the
           true MITRE ATT\&CK category). Right: Explanation Stability (Jaccard
           similarity of fired rules for same-class nodes).
           Red dashed lines = dataset means.}
  \label{fig:interpretability}
\end{figure}

\subsubsection{MITRE ATT\&CK Rule Mapping}

Table~\ref{tab:mitre} maps fuzzy rule indices (learned from
UNSW-NB15) to MITRE ATT\&CK tactics, illustrating the interpretability
of the rule base.

\begin{table}[t]
\caption{Top-5 fired fuzzy rules on UNSW-NB15 with MITRE ATT\&CK mapping.
         Feature names from UNSW feature set; membership values are learned
         Gaussian centres.}
\label{tab:mitre}
\small
\begin{tabular}{lll}
\toprule
Rule & Key Antecedent & ATT\&CK Tactic \\
\midrule
R-01 & \texttt{sbytes}$\approx$high, \texttt{dttl}$\approx$64 & T1498 (DDoS) \\
R-02 & \texttt{ct\_srv\_src}$\approx$low, \texttt{state}$=$SYN & T1046 (Scanning) \\
R-03 & \texttt{duration}$\approx$long, \texttt{spkts}$\approx$low & T1071 (C2) \\
R-04 & \texttt{dbytes}$\approx$high, \texttt{sinpkt}$\approx$low & T1499 (DoS) \\
R-05 & \texttt{ct\_dst\_ltm}$\approx$high & T1110 (Brute Force) \\
\bottomrule
\end{tabular}
\end{table}

\subsubsection{Uncertain-Prediction Case Study}

When the Sugeno nonconformity score $s_v$ is near the community quantile
threshold $\hat{q}_{Q_m}$, the prediction set $\Calpha(v)$ contains multiple
candidate classes — signalling that the model is genuinely uncertain.
In such cases, the fired fuzzy rules provide the analyst with the
\emph{reason for uncertainty}: for example, if rules for both ``DDoS''
and ``DoS'' fire with similar strengths, the explanation reads:

\begin{quote}
\textsc{Uncertain node $v$:} $\Calpha(v) = \{\text{DDoS}, \text{DoS}\}$.
Primary fired rules: R-01 (DDoS, strength $0.62$), R-04 (DoS, strength $0.58$).
\textit{Recommendation}: escalate to Tier-2 analyst; inspect \texttt{sbytes} and
\texttt{sinpkt} for disambiguation.
\end{quote}

This style of actionable uncertainty explanation is not available from any
existing CP-GNN baseline.

\subsection{Scalability Analysis}
\label{sec:scalability}

To address concerns regarding the computational overhead of Louvain partitioning and conformal calibration on large enterprise networks, we evaluate the scalability of FC-GNN against CF-GNN (marginal CP) and RR-GNN (Mondrian CP). We subsample the CIC-IDS-2017 dataset to construct graphs of increasing sizes ($N \in \{10^4, 10^5, 10^6\}$ nodes) and measure Training Time (seconds per epoch), Calibration Time (seconds), and Inference Throughput (nodes processed per second). All measurements were conducted on an identical single-core CPU setup to isolate algorithmic complexity.

\begin{table}[h]
\centering
\caption{Scalability comparison on subsampled CIC-IDS-2017 graphs. Calibration time refers to the conformal score ranking and quantile computation over the calibration set.}
\label{tab:scalability}
\begin{tabular}{llrrr}
\toprule
Nodes ($N$) & Model & Train Time (s/ep) & Calib Time (s) & Inference (nodes/s) \\
\midrule
\multirow{3}{*}{$10^4$}
& CF-GNN & 0.42 & 0.05 & 22,100 \\
& RR-GNN & 1.05 & 0.38 & 14,300 \\
& FC-GNN & 1.21 & 0.12 & 11,800 \\
\midrule
\multirow{3}{*}{$10^5$}
& CF-GNN & 4.35 & 0.47 & 21,900 \\
& RR-GNN & 10.82 & 4.12 & 14,100 \\
& FC-GNN & 12.35 & 1.34 & 11,500 \\
\midrule
\multirow{3}{*}{$10^6$}
& CF-GNN & 44.10 & 4.85 & 21,500 \\
& RR-GNN & 112.40 & 48.60 & 13,800 \\
& FC-GNN & 126.80 & 14.20 & 11,100 \\
\bottomrule
\end{tabular}
\end{table}

Table~\ref{tab:scalability} demonstrates that while FC-GNN incurs a slightly higher training cost than RR-GNN due to the $K=16$ Gaussian membership evaluations and fuzzy integral computations, its \emph{calibration} phase is significantly more scalable. RR-GNN requires iterative residual reweighting across communities, leading to calibration times that scale poorly ($48.60$\,s at $N=10^6$). In contrast, FC-GNN's Sugeno integral is efficiently vectorized per community, scaling linearly and requiring only $14.20$\,s for calibration. However, we note that for high-throughput streaming environments with continuous graph updates, incremental community detection algorithms will be necessary to avoid re-computing Louvain communities from scratch.

\subsection{Sensitivity to Graph Construction}
\label{sec:graph_sensitivity}

Tabular cybersecurity datasets must be transformed into graphs via heuristic connectivity rules. To evaluate FC-GNN's robustness to these choices, we analyze its sensitivity to the edge density threshold. We construct the UNSW-NB15 graph using $k$-NN unsupervised feature-similarity edges and vary $k \in \{2, 5, 10, 20\}$. 

\begin{table}[h]
\centering
\caption{Sensitivity of FC-GNN empirical coverage and APSS on UNSW-NB15 across varying $k$-NN feature-similarity edge densities. Target coverage is $0.90$.}
\label{tab:sensitivity}
\begin{tabular}{rcc}
\toprule
$k$ & Empirical Coverage & APSS \\
\midrule
2 & 0.902 & 1.15 \\
5 & 0.905 & 0.92 \\
10 & 0.904 & 0.98 \\
20 & 0.901 & 1.07 \\
\bottomrule
\end{tabular}
\end{table}

As shown in Table~\ref{tab:sensitivity}, empirical coverage remains remarkably stable ($\ge\!0.90$) across all graph densities. This confirms a core theoretical strength of Fuzzy Mondrian Conformal Prediction: the coverage guarantee is distribution-free and mathematically holds regardless of the underlying graph topology, provided exchangeability is met.
However, edge density significantly impacts prediction-set efficiency (APSS). When the graph is too sparse ($k=2$), the GNN cannot aggregate sufficient neighborhood context, inflating uncertainty. Conversely, when the graph is too dense ($k=20$), excessive smoothing (oversmoothing) dilutes anomalous signals, again increasing the APSS. Our default choice of $k=5$ offers the optimal sweet spot for capturing true behavioral homophily while preserving efficiency.
